\subsection{Limitations of existing systems and research gaps}
\label{sec:lr-limitations}

\subsubsection{Imperfect substrates demand architectures that expose, not bury, their flaws}

No open scholarly knowledge graph is a ground truth, and the architecture
chosen to consume one determines whether its imperfections remain visible
or are absorbed into latent representations.  Citation recall across the
major indexes varies by a factor of three, from 88\% in Google Scholar to
28\% in OpenCitations' COCI, with systematic disciplinary bias across
providers \cite{martinmartin2021coverage}.  OpenAlex itself, the substrate
this thesis ingests, retains references to deleted works and does not
systematically index non-source references \cite{culbert2025reference},
and over sixty percent of its journal articles lack complete institutional
metadata, with Chinese institutions particularly under-represented
\cite{zhang2024missing}.  Single-vector dense retrieval and opaque RAG
pipelines compound these gaps by collapsing them into representations the
user cannot interrogate.  A KG-native architecture instead surfaces
missing affiliations, contested author identities and shifting topic
labels as queryable structure, making them features of retrieval rather
than latent corruption. \\

\subsubsection{No method, system or benchmark addresses the joint task this thesis targets}
\label{sec:lr-gap}

Each thread of prior work identifies a single-axis gap, and the thesis's
contribution lies in their unaddressed intersection.  Methodologically,
BM25 is lexically brittle, dense retrievers miss roughly a quarter of
relevant scientific results (Section~\ref{sec:lr-dense}), citation-graph
methods are conceptually well-matched but unbenchmarked on
natural-language scholarly queries (Section~\ref{sec:lr-graph}),
structured-query translation has been evaluated only on
template-generated scholarly graphs (Section~\ref{sec:lr-nl2q}), and
keyphrase extraction is scored by agreement with author gold sets rather
than by the retrieval an index built from it supports
(Section~\ref{sec:lr-keyphrase}).  At the
system level, deployed engines split cleanly into text-centric tools
that treat the citation graph as a ranking signal at most and
graph-native discovery tools that abandon natural-language entry
entirely.  The LLM-augmented systems consume both as black boxes
(Sections~\ref{subsec:traditional-search}
and~\ref{subsec:llm-agentic}).
At the benchmark level, no test set jointly exercises natural-language
scientific queries, structured-query translation, and graph traversal
over a live scholarly substrate.  These three gaps converge on a single
absence.  No current system, on any current benchmark, with any
documented method combination, performs the joint task that a KG-native
hybrid scholarly retrieval pipeline would. \\

\subsubsection{The thesis's contribution}

This thesis addresses that absence with a system that uses an LLM to
translate natural-language questions into a structured intermediate
representation, compiles that representation deterministically into
optimised Cypher over an OpenAlex-derived schema, and arranges BM25,
dense, and graph-traversal evidence as successive stages of a single
retrieval cascade, in which the structured stage gates the candidate set
and the dense and cross-encoder stages rank and re-rank it.  
Architecturally, this treats the knowledge graph as the
primary semantic interface rather than an auxiliary feature store, with
BM25 and dense retrieval acting as complementary lexical and semantic
lenses on the same graph-anchored universe of works, authors, venues and
citations.  The agentic and RAG systems surveyed in
Section~\ref{subsec:llm-agentic} are users of such an architecture rather
than competitors to it, consuming scholarly retrieval as a tool rather
than providing it as a substrate.